\section{Introduction}
\label{sec:intro}

This report documents the team assessment of the Database Management Systems (DBMS) course.
Using MongoDB as the DBMS assigned by the lecturer and two public datasets, we study how a
real system supports \emph{indexing}, \emph{query processing and optimization},
\emph{transactions} and \emph{concurrency control}, and we systematically compare the theory
taught in class with the technology used in practice.

\subsection{Objectives and case studies}
\label{sec:objectives}

The assessment asks the team to show how the DBMS works on sample data for the following four
topics, and to explain each one against the course theory.
\begin{enumerate}
  \item How the DBMS works with \textbf{indexing} (\Cref{sec:indexing}).
  \item How the DBMS works with \textbf{query processing and optimization}
        (\Cref{sec:query}).
  \item How the DBMS works with \textbf{transactions} (\Cref{sec:transaction}).
  \item How the DBMS works with \textbf{concurrency control} (\Cref{sec:concurrency}).
\end{enumerate}
All topics are demonstrated on one common application with three required case studies:
\begin{enumerate}[label=\textbf{C\arabic*.}]
  \item retrieve the text record and the images of a patient given a \emph{patient ID};
  \item retrieve the text record and the images given a patient ID \emph{and conditions}
        (for example, only T2 sagittal MRI images);
  \item retrieve the MRI images that are \emph{similar} to one MRI image of a given patient.
\end{enumerate}

\subsection{Datasets}
\label{sec:datasets}
The application uses two public datasets from Mendeley Data: a text dataset with the patients'
records and an image dataset with their MRI scans. \todo{Describe each dataset: number of
patients and files, fields, image types and sizes, any cleaning or additional data we added, and
the final schema loaded into MongoDB.}

\subsection{Organization of the report}
\label{sec:organization}
\Cref{sec:mongodb-overview} gives an overview of MongoDB. \Cref{sec:indexing} to
\Cref{sec:concurrency} treat the four topics; each of them is organized in the same way: the
\emph{theory in class}, the \emph{implementation in MongoDB}, a \emph{demonstration} on the
case studies, and a \emph{comparison} between theory and practice. \Cref{sec:conclusion}
concludes and \Cref{sec:contrib} lists the contribution of each member.

\subsection{MongoDB Overview}
\label{sec:mongodb-overview}

MongoDB is an open, general-purpose \emph{document-oriented} NoSQL database
\cite{mongodb-intro}. Rather than spreading an entity over several fixed-schema tables that must
be joined at query time, it stores each entity as one self-describing \emph{document}. This suits
semi-structured data such as patient records, whose attributes vary from one patient to another,
and it lets the database scale out over many machines.

\subsubsection{Data model}
A MongoDB deployment hosts \emph{databases}; a database holds \emph{collections}; a collection
holds \emph{documents}. \Cref{tab:terms} maps relational terms to MongoDB terms.

\begin{table}[H]
\centering
\begin{tabularx}{0.85\textwidth}{@{}>{\raggedright\arraybackslash}X>{\raggedright\arraybackslash}X@{}}
\toprule
\textbf{Relational DBMS} & \textbf{MongoDB} \\
\midrule
Database & Database \\
Table & Collection \\
Row (tuple) & Document \\
Column (attribute) & Field \\
Primary key & \code{\_id} field \\
Join & Embedded documents, \code{\$lookup} stage \\
Schema fixed per table & Flexible per document (optional validation rules) \\
\bottomrule
\end{tabularx}
\caption{Relational terminology versus MongoDB terminology.}
\label{tab:terms}
\end{table}

\paragraph{Documents and BSON.}
A document is an ordered set of field--value pairs. Internally and on the wire, documents are
encoded in \emph{BSON}, a binary superset of JSON. BSON adds types that JSON lacks, such as
\code{Int32}, \code{Int64}, \code{Double}, \code{Decimal128}, \code{Date}, \code{ObjectId} and
\code{Binary}, and it is designed to be fast to traverse and to encode. A value can itself be a
nested document or an array, so a hierarchy of related data can live in one document. A single
document may not exceed 16~MB, which is why large binary objects such as images are normally kept
outside the document (on disk, or in GridFS, which splits a file into chunks) while the document
stores only metadata and a path.

\paragraph{The \code{\_id} field.}
Every document has a unique \code{\_id}, which acts as its primary key. If the client does not
supply one, the driver generates a 12-byte \code{ObjectId} made of a 4-byte timestamp, a 5-byte
random value and a 3-byte counter, so identifiers are practically unique and roughly increasing
in time without any coordination between clients. MongoDB automatically builds a unique index on
\code{\_id} for every collection.

\paragraph{Flexible schema and schema validation.}
By default, documents in the same collection may have different fields and field types. When
stricter control is needed, a collection can carry a JSON Schema validator that rejects documents
which do not conform. Flexibility makes it easy to evolve the data over time, but it moves part of
the responsibility for data quality from the DBMS to the application.

\paragraph{Embedding versus referencing.}
Modeling in MongoDB is driven by the access pattern rather than by normalization.
\Cref{tab:embed} contrasts the two basic designs. In our application, the case studies always
start from a patient ID and then need the patient's text record together with images, so a
patient and its image metadata are read together; at the same time, images are numerous and are
also queried on their own (by series and orientation, and by similarity). \todo{State the final
design chosen by the team: one collection per entity with references, or images embedded in
the patient document.}

\begin{table}[H]
\centering
\small
\begin{tabularx}{\textwidth}{@{}>{\raggedright\arraybackslash}p{2.6cm}>{\raggedright\arraybackslash}X>{\raggedright\arraybackslash}X@{}}
\toprule
 & \textbf{Embedding} & \textbf{Referencing} \\
\midrule
Idea & Nest related data inside one document & Keep related data in separate documents linked by an identifier (such as \code{patient\_id}) \\
Reads & One read returns everything; no join & Needs a second query or \code{\$lookup} \\
Writes & A document grows; atomic as one document & Each document is small; updates touch several documents \\
Best when & Data is read together and its size is bounded & Data is large, unbounded, shared, or queried independently \\
\bottomrule
\end{tabularx}
\caption{Embedding versus referencing in MongoDB data modeling.}
\label{tab:embed}
\end{table}

\begin{lstlisting}[language=mongosh,caption={A patient document that embeds image metadata}]
{
  _id: ObjectId("..."),
  patient_id: "P001",
  age: 54,
  report: "Mild disc bulge at L4-L5 ...",
  images: [
    { series: "T2", orientation: "sagittal", path: "img/P001_t2_sag_01.png" },
    { series: "T1", orientation: "axial",    path: "img/P001_t1_ax_01.png" }
  ]
}
\end{lstlisting}

\subsubsection{Query interface}
Applications talk to MongoDB through official drivers (for example PyMongo for Python) or the
interactive shell \code{mongosh}. CRUD is expressed by methods such as \code{insertOne},
\code{insertMany}, \code{find}, \code{updateOne} and \code{deleteOne}. A query is itself a
document: field conditions are combined by default with AND, and operators such as \code{\$eq},
\code{\$gt}, \code{\$in}, \code{\$and} and \code{\$or} build richer conditions. Projection chooses
the returned fields, and \code{sort} and \code{limit} shape the result. Complex analysis uses the
\emph{aggregation pipeline}, a sequence of stages (\code{\$match}, \code{\$group},
\code{\$sort}, \code{\$lookup}, \code{\$project}, \dots) where each stage transforms the stream of
documents produced by the previous one. The three case studies map to queries of the following
form.

\begin{lstlisting}[language=mongosh,caption={Query forms of the three case studies (sketch)}]
// C1: by patient ID
db.images.find({ patient_id: "P001" })

// C2: by patient ID and conditions
db.images.find({ patient_id: "P001", series: "T2", orientation: "sagittal" })

// C3: similar images (vector search stage)
db.images.aggregate([
  { $vectorSearch: { index: "image_vec_idx", path: "embedding",
                     queryVector: queryEmbedding, numCandidates: 200, limit: 5 } }
])
\end{lstlisting}

\subsubsection{Architecture}
\label{sec:mongodb-arch}
The database server process is \code{mongod}. It receives requests from clients, executes them
through the query engine, and delegates storage to a pluggable \emph{storage engine}.

\paragraph{Storage engine: WiredTiger.}
The default storage engine is WiredTiger \cite{mongodb-storage}. It keeps collections and indexes
in B-tree-based structures on disk and caches the working set in memory (by default about half of
the RAM minus 1~GB). It provides \emph{document-level concurrency control}: concurrent writers lock
individual documents rather than whole collections, and readers see a consistent snapshot of the
data (multiversion concurrency control). Durability is obtained by combining a write-ahead
\emph{journal} with periodic \emph{checkpoints} of the data files (every 60 seconds by default).
Data is compressed on disk (Snappy by default for collections) and index keys use prefix
compression, which reduces I/O and storage.

\paragraph{Replication.}
A \emph{replica set} is a group of \code{mongod} instances holding the same data. One member is the
\emph{primary} and receives all writes; the \emph{secondaries} replicate the primary's operation log
(\emph{oplog}) and apply it. If the primary becomes unavailable, the remaining members hold an
election and promote a secondary automatically, giving redundancy and high availability
\cite{mongodb-replication}. Write concern and read concern let the application choose how many
members must acknowledge a write and how up to date a read must be.

\paragraph{Sharding.}
To scale beyond one machine, a collection can be partitioned horizontally into \emph{chunks} spread
over several \emph{shards} (each usually a replica set) according to a \emph{shard key}. A router
process (\code{mongos}) directs each operation to the relevant shards using metadata kept by the
config servers. The choice between a ranged and a hashed shard key trades efficient range queries
against an even distribution of data.

\paragraph{Transactions.}
Every write to a single document, including its nested fields and arrays, is atomic. For
operations spanning several documents, MongoDB supports multi-document ACID transactions on
replica sets (since version~4.0) and on sharded clusters (since version~4.2), using snapshot
isolation. Transactions are treated in \Cref{sec:transaction}.

\paragraph{Indexes.}
MongoDB supports single-field, compound, multikey, hashed, text, geospatial, wildcard, TTL and
partial indexes, among others, which are examined in \Cref{sec:indexing}.

\paragraph{Search and vector search.}
Full-text search and vector search are provided by a separate process, \code{mongot}, that is
deployed beside \code{mongod} and keeps its own search indexes synchronized with the data. They
became generally available for the self-managed Community Edition with MongoDB~8.2 or later
\cite{mongodb-search-ga}. Vector search is the basis of our similar-image retrieval (case study
C3).

\subsubsection{Editions and deployment}
MongoDB is offered as the Community Edition (free, source-available), the commercial Enterprise
edition for self-managed use, and the cloud service Atlas. This report uses the self-managed
Community Edition. \todo{Give the exact MongoDB version, the operating system, how it was installed
(for example with Docker), and the hardware used for the experiments.}

\subsubsection{Strengths and limitations}
\begin{description}
  \item[Strengths.] A flexible schema suited to evolving, semi-structured data; a data model
        close to application objects; rich indexing; built-in replication and horizontal scaling;
        integrated search and vector search.
  \item[Limitations.] Joins across collections are less efficient than in relational systems;
        denormalized (embedded) data may be duplicated and must be kept consistent by the
        application; documents are limited to 16~MB; the working set should fit in memory for the
        best performance.
\end{description}
